% 狮子座（综述）
% license CCBYSA3
% type Wiki

本文根据 CC-BY-SA 协议转载翻译自维基百科\href{https://en.wikipedia.org/wiki/Leo_(constellation)}{相关文章}。

狮子座（Leo，/ˈliːoʊ/）是黄道十二星座之一，位于西侧的巨蟹座（Cancer，螃蟹）与东侧的室女座（Virgo，处女）之间，坐落在北天球。其名称源自拉丁语，意为“狮子”，在古希腊文化中代表神话英雄赫拉克勒斯（Heracles）在其十二项功绩之一中所杀死的尼米亚狮子。其古老的天文符号为♌︎。作为公元2世纪天文学家托勒密所描述的48个星座之一，狮子座至今仍是现代88个星座之一，并且由于其众多明亮的恒星以及一个极具辨识度、令人联想到伏地狮子的独特形状，狮子座也是最容易识别的星座之一。
\subsection{特征}
\subsubsection{恒星}
\begin{figure}[ht]
\centering
\includegraphics[width=6cm]{./figures/cd1a625b724681d2.png}
\caption{狮子座的肉眼可见外观（图像中央的明亮天体为2004年3月的行星木星）。} \label{fig_Leo_1}
\end{figure}
狮子座包含许多明亮恒星，其中许多在古代就已被单独识别。该星座中有九颗明亮恒星可以轻易用肉眼观测到，其中四颗属于一等或二等星，这使得狮子座显得尤为醒目。这九颗恒星中的六颗还组成了一个称为“镰刀”（The Sickle）的星群，对于现代观测者来说，其形状类似一个反向的“问号”。镰刀由六颗恒星标记：狮子座$\varepsilon$、狮子座$\mu$、狮子座$\zeta$、狮子座$\gamma$、狮子座$\eta$以及狮子座$\alpha$。其余三颗恒星构成一个等腰三角形，其中狮子座$\beta$（Denebola）标示狮子的尾巴，而狮子的其余身体部分则由狮子座$\delta$和狮子座$\theta$勾勒出来$^{[2]}$。
\begin{itemize}
\item 狮子座$\alpha$，称为轩辕十四（Regulus），是一颗蓝白色主序星，视星等为1.34，距离地球77.5光年。它是一对双星，用双筒望远镜即可分辨，其伴星的视星等为7.7。其传统名称Regulus意为“小国王”。
\item 狮子座$\beta$，又称Denebola，位于星座中与轩辕十四相对的一端。它是一颗蓝白色恒星，视星等为2.23，距离地球36光年。Denebola一名意为“狮子的尾巴”。
\item 狮子座$\gamma$，又称Algieba，是一个双星系统，并伴有一颗视线上的第三成员；主星和伴星可用小型望远镜分辨，第三成员可用双筒望远镜看到。主星是一颗金黄色巨星，视星等为2.61，伴星性质相似但视星等为3.6；它们的轨道周期约为600年，距离地球126光年。无关的第三星，即狮子座40，是一颗略带黄色的恒星，视星等为4.8。其传统名称Algieba意为“前额”。
\item 狮子座$\delta$，称为Zosma，是一颗蓝白色恒星，视星等为2.58，距离地球58光年。
\item 狮子座$\mu$，称为Rasalas，是一颗红色恒星，视星等为3.88，亮度足以用肉眼观测。该系统距离地球124光年。Rasalas亦称Alshemali，这两个名称均源自“Al Ras al Asad al Shamaliyy”的缩写，意为“狮子朝南的头部”。
\item 狮子座$\theta$，称为Chertan，是一颗白色恒星，视星等为3.324，可用肉眼观测，是该星座中较亮的恒星之一，距离地球约165光年。Chertan一名源自阿拉伯语al-kharātān，意为“两根小肋骨”，最初指的是狮子座$\delta$和狮子座$\theta$，$^{[3]}$。
\end{itemize}
\begin{figure}[ht]
\centering
\includegraphics[width=8cm]{./figures/a1a3d9c69259c7b9.png}
\caption{狮子座最亮的恒星} \label{fig_Leo_2}
\end{figure}
\begin{figure}[ht]
\centering
\includegraphics[width=6cm]{./figures/fe6607081f71245d.png}
\caption{狮子座，上方是小狮子座，如乌拉尼亚之镜所描绘，这是一套约1825年在伦敦出版的星座卡片} \label{fig_Leo_3}
\end{figure}
狮子座中还包含一颗明亮的变星——红巨星狮子座R（R Leonis）。它是一颗米拉型变星，最暗时视星等约为10，通常的最亮视星等为6；并且会周期性地增亮至视星等4.4。狮子座R距离地球约330光年，周期为310天，直径约为太阳直径的450倍$^{[4]}$。

恒星沃尔夫359（Wolf 359，又名狮子座CN）位于狮子座中，是距离地球最近的恒星之一，距离仅7.8光年。沃尔夫359是一颗红矮星，视星等为13.5；由于它是一颗耀斑星，其亮度会周期性地增亮不超过一个星等$^{[4]}$。格利泽436（Gliese 436）是一颗位于狮子座内、距离太阳约33光年的暗弱恒星，其周围运行着一颗凌日的海王星质量级系外行星$^{[5]}$。

碳星CW Leo（IRC +10216）是在红外N波段（波长10\,$\mu$m）夜空中最明亮的恒星。

恒星SDSS J102915+172927（又称Caffau之星）是一颗位于狮子座方向的银河系晕族II恒星，其年龄约为130亿年，是银河系中已知最古老的恒星之一。它具有目前已知恒星中最低的金属丰度。

现代天文学家（包括1602年的第谷·布拉赫）将一组原本构成“狮子尾端毛簇”的恒星从狮子座中剔除，并用它们建立了新的星座——后发座（Berenice之发），尽管在古希腊和古罗马时期，这一名称已有先例$^{[6]}$。
\subsubsection{深空天体}
狮子座中包含许多明亮的星系，其中最著名的包括梅西耶65、梅西耶66、梅西耶95、梅西耶96、梅西耶105以及NGC\,3628$^{[7]}$，其中前两者属于著名的“狮子座三重星系”$^{[8]}$。

狮子座环（Leo Ring）是一团由氢和氦气体组成的云，位于该星座内两座星系的轨道附近。
\begin{figure}[ht]
\centering
\includegraphics[width=6cm]{./figures/3b150d76df662684.png}
\caption{梅西耶\,66} \label{fig_Leo_4}
\end{figure}
M66 是一座螺旋星系，属于狮子座三重星系，其另外两个成员是 M65 和 NGC\,3628。它距离地球约 3700 万光年，由于与三重星系中其他成员之间的引力相互作用，其形态略显扭曲，这些相互作用正在将恒星从 M66 中拉离。最终，最外层的恒星可能会形成一座绕 M66 运行的矮星系$^{[9]}$。M65 和 M66 均可通过大型双筒望远镜或小型望远镜观测到，但它们高度集中的核区以及拉伸的形态特征，只有在大型业余观测设备中才能清晰辨认$^{[4]}$。
\begin{figure}[ht]
\centering
\includegraphics[width=6cm]{./figures/04a0b3e20a43eeff.png}
\caption{著名的引力透镜——被称为“宇宙马蹄”（Cosmic Horseshoe）} \label{fig_Leo_5}
\end{figure}
M95 和 M96 均为距离地球约 2000 万光年的螺旋星系。尽管在小型望远镜中它们呈现为模糊的光斑，但其结构特征只有在较大的观测设备中才能显现。M95 是一座棒旋星系。M105 位于 M95/M96 星系对约一度之外，是一座视星等约为 9 等的椭圆星系，同样距离地球约 2000 万光年$^{[4]}$。

NGC\,2903 是一座由威廉·赫歇尔于 1784 年发现的棒旋星系，其大小和形态与银河系非常相似，距离地球约 2500 万光年。在其核心区域，NGC\,2903 拥有许多“热点”，这些热点被发现位于恒星形成区域附近。该区域的恒星形成被认为与星系中富含尘埃的棒结构有关，该结构在旋转过程中向直径约 2000 光年的区域传递冲击波。星系的外缘分布着大量年轻的疏散星团$^{[9]}$。

狮子座还包含了一些可观测宇宙中最大的结构。其中包括 Clowes–Campusano LQG、U1.11、U1.54 以及 Huge-LQG，它们均为大型类星体群；其中 Huge-LQG 是目前已知的第二大结构$^{[10]}$（另见 NQ2–NQ4 伽马射线暴过密区）。
\subsubsection{流星雨}
狮子座流星雨（Leonids）出现在每年 11 月，峰值通常在 11 月 14–15 日，其辐射点靠近狮子座 γ 星。其母体为坦普尔–塔特尔彗星，该彗星每隔约 35 年会引发显著的流星暴。通常的峰值流量约为每小时 10 颗流星$^{[2]}$。

一月狮子座流星雨（January Leonids）是一场较小的流星雨，活动高峰出现在 1 月 1 日至 7 日之间$^{[11]}$。
\subsection{历史与神话}
\begin{figure}[ht]
\centering
\includegraphics[width=6cm]{./figures/9c7c060f8627f118.png}
\caption{约公元\,1000\,年的一份西方科学手稿中的狮子座} \label{fig_Leo_6}
\end{figure}
狮子座是最早被识别的星座之一，考古证据表明，美索不达米亚人在公元前约 4000 年就已经拥有一个与之相似的星座$^{[12]}$。波斯人称狮子座为 Ser 或 Shir；土耳其人称其为 Artan；叙利亚人称其为 Aryo；犹太人称其为 Arye；印度人称其为 Simha，这些名称的含义均为“狮子”。
\begin{figure}[ht]
\centering
\includegraphics[width=6cm]{./figures/168668ef981f00db.png}
\caption{塞浦路斯邮票，描绘了赫拉克勒斯与狮子相遇的马赛克图像——涅墨亚狮} \label{fig_Leo_7}
\end{figure}
一些神话学家认为，在苏美尔文化中，狮子座代表怪物洪巴巴（Humbaba），它被吉尔伽美什所杀$^{[13]}$。

在巴比伦天文学中，该星座被称为 UR.GU.LA，即“伟大的狮子”；其亮星轩辕十四（Regulus）被称为“位于狮子胸前的恒星”。轩辕十四还具有鲜明的王权象征意义，因为它被称为“王者之星”$^{[14]}$。

在希腊神话中，狮子座被认定为涅墨亚狮（Nemean Lion），它是在赫拉克勒斯（罗马人称其为赫丘利）完成十二项伟业中的第一项时被杀死的$^{[12][2]}$。涅墨亚狮会将女子劫持到洞穴巢穴中，引诱附近城镇的勇士前来营救受困少女，而这些勇士往往因此丧命$^{[15]}$。这头狮子刀枪不入，因此战士们的棍棒、宝剑和长矛对它毫无作用。赫拉克勒斯意识到必须徒手制服狮子，于是潜入其洞穴，与之展开近身搏斗$^{[15]}$。当狮子扑来时，赫拉克勒斯在半空中将其抓住，一只手握住狮子的前腿，另一只手抓住后腿，将其向后弯折，折断脊背，从而解救了被囚禁的少女$^{[15]}$。宙斯为了纪念这一伟业，将狮子安置于天空之中$^{[15]}$。

罗马诗人奥维德称其为 Herculeus Leo（赫拉克勒斯之狮）和 Violentus Leo（狂暴之狮）。Bacchi Sidus（酒神之星）也是它的称号之一，因为酒神巴克科斯常与狮子联系在一起。然而，曼尼利乌斯则称其为 Jovis et Junonis Sidus（朱庇特与朱诺之星）。
\subsection{占星学}  
截至 2002 年，太阳在天文学意义上的狮子座内运行的时间为 8 月 10 日至 9 月 16 日$^{[16]}$。在回归占星学中，太阳被认为在 7 月 23 日至 8 月 22 日位于狮子座；在恒星占星学中，则为 8 月 16 日至 9 月 17 日。
\subsection{同名事物}  
USS Leonis（AK-128）是一艘美国海军陨石坑级货船。
\subsection{另见} 
\begin{itemize}
\item 狮子座（中国天文学）
\end{itemize}
\subsection{参考文献}
\begin{enumerate}
\item 柯克帕特里克，J·戴维；马罗科，费德里科等（2024年4月）。“基于全天空20秒差距范围内约3600颗恒星和棕矮星普查的初始质量函数”. 天体物理学杂志增刊系列. 271 (2): 55. arXiv:2312.03639. 文献标识码:2024ApJS..271...55K. doi:10.3847/1538-4365/ad24e2.
\item 里德帕斯与蒂里昂 2001, 第166–167页。
\item “国际天文学联合会批准的恒星名称列表”。于2025年3月10日从原始网页存档。检索日期：2022年1月24日。
\item 里德帕斯与蒂里昂 2001, 第166–168页。
\item 天文学家发现迄今最小的“系外行星”。多伦多。2009年1月16日存档自原文。
\item L. 菲尔·辛普森（施普林格，2012）星座指南：望远镜观测、传说与神话，第235页（ISBN 9781441969415).
\item 詹姆斯·米勒。“狮子座的深空天体”。天文之旅。检索日期：2025年12月21日：CS1维护：网址状态（链接)
\item 安托万；达莉娅·格雷林（2023年6月5日）。“狮子座三重星系 | M65、M66 和 NGC 3628 天体摄影”。银河猎手。检索日期：2025年12月21日：CS1维护：网址状态（链接)
\item 威尔金斯，杰米；邓恩，罗伯特（2006）。300个天文天体：宇宙视觉参考。纽约布法罗：萤火虫图书。ISBN 978-1-55407-175-3.
\item 普罗斯塔克，塞尔吉奥（2013年1月11日）。“发现宇宙中最大的结构”。scinews.com。检索日期：2013年1月15日。
\item 詹尼斯肯斯，彼得（2011年9月）。“绘制流星体轨道：新流星雨被发现”。天空与望远镜: 24.
\item 帕萨乔夫，杰伊·M.（2006）。恒星与行星。马萨诸塞州波士顿：霍顿·米夫林。ISBN 9780395537596.
\item 塔姆拉·安德鲁斯（牛津大学出版社，2000年）自然神话词典：大地、海洋与天空的传说（ISBN 9780195136777).
\item 巴比伦星象学加文·怀特著，索拉里亚出版社，2008年，第140页ISBN 978-0955903700
\item 珍妮特·帕克等编（2007）。神话：神话、传说与幻想。斯特鲁克。第121–122页。ISBN9781770074538
\item 让·梅厄斯，《黄道星座》，载于：更多数学天文学小品（里士满：威尔曼-贝尔公司，2002年），第327-333页。其他年份的日期可能因插入闰日.
\end{enumerate}
\begin{itemize}
\item 《恒星名称：其传说与含义》，理查德·艾伦·欣克利著，多佛出版社出版。ISBN 0-486-21079-0
\item 里德帕斯，伊恩；蒂里昂，威尔（2001）恒星与行星指南, 普林斯顿大学出版社,ISBN 0-691-08913-2
\item 伊恩·里德帕斯和威尔·蒂里昂（2007）。恒星与行星指南, 科林斯，伦敦。ISBN 978-0-00-725120-9。普林斯顿大学出版社，普林斯顿。ISBN 978-0-691-13556-4.
\item 符号词典, 卡尔·G·利翁曼著，W. W. 诺顿公司。ISBN 0-393-31236-4
\end{itemize}
\subsection{外部链接}
\begin{itemize}
\item 可点击的狮子座
\item 瓦尔堡研究所图像数据库（关于利奥的中世纪与早期现代图像）
\item 来自StarDate Online的信息 已归档2023-06-05于互联网档案馆
\item 伊恩·里德帕斯的星图——狮子座
\end{itemize}